\documentclass[12pt]{article}

% Packages
\usepackage{amsmath, amssymb, amsthm}
\usepackage[margin=1in]{geometry}
\usepackage{hyperref}

% Theorem environments — separate counters
\newtheorem{theorem}{Theorem}
\newtheorem{lemma}{Lemma}
\newtheorem{corollary}{Corollary}
\newtheorem{conjecture}{Conjecture}
\theoremstyle{definition}
\newtheorem{definition}{Definition}
\theoremstyle{remark}
\newtheorem{remark}{Remark}

% Title
\title{On the Greatest Common Divisor of Consecutive Catalan Numbers}
\author{Daniel Brooks \\ \small University of Chicago}
\date{\today}

\begin{document}

\maketitle

\begin{abstract}
We investigate the greatest common divisor of consecutive Catalan numbers $C_n$ and $C_{n+1}$, and establish conditions under which consecutive Catalan numbers share a common factor.
\end{abstract}

\section{Introduction}

The Catalan numbers are a sequence of natural numbers defined by
\[
C_n = \frac{(2n)!}{(n+1)!\, n!}
\]
for $n \geq 0$. The first several terms are
\[
1, 1, 2, 5, 14, 42, 132, 429, 1430, 4862, \ldots
\]

These numbers arise in numerous combinatorial contexts, including the enumeration of valid parenthesizations, full binary trees, and triangulations of polygons.

A well-known recurrence relation for the Catalan numbers is
\begin{equation}\label{eq:recurrence}
C_{n+1} = \frac{2(2n+1)}{n+2} \cdot C_n.
\end{equation}

In this paper, we study the quantity $\gcd(C_n, C_{n+1})$ and seek to characterize when consecutive Catalan numbers are coprime and when they share a common factor.

\newpage
\section{Preliminary Results}

We compute $\gcd(C_n, C_{n+1})$ for small values of $n$:
\begin{center}
\begin{tabular}{c|c|c|c}
$n$ & $C_n$ & $C_{n+1}$ & $\gcd(C_n, C_{n+1})$ \\
\hline
0 & 1 & 1 & 1 \\
1 & 1 & 2 & 1 \\
2 & 2 & 5 & 1 \\
3 & 5 & 14 & 1 \\
4 & 14 & 42 & 14 \\
5 & 42 & 132 & 6 \\
6 & 132 & 429 & 33 \\
7 & 429 & 1430 & 143 \\
8 & 1430 & 4862 & 286 \\
\end{tabular}
\end{center}

From Equation~\eqref{eq:recurrence}, we note that $C_{n+1}$ is determined by multiplying $C_n$ by the fraction $\frac{2(2n+1)}{n+2}$. The key quantity to analyze is therefore $\frac{2(2n+1)}{n+2}$ and its relationship with $C_n$.

\section{Main Results}

\begin{theorem}\label{thm:divisibility}
$C_n \mid C_{n+1}$ if and only if $n \in \{0, 1, 4\}$.
\end{theorem}

\begin{proof}
From the recurrence relation~\eqref{eq:recurrence}, $C_n \mid C_{n+1}$ if and only if $\frac{2(2n+1)}{n+2}$ is a positive integer.

We rewrite the numerator:
\[
2(2n+1) = 4n + 2 = 4(n+2) - 6.
\]
Therefore,
\[
\frac{2(2n+1)}{n+2} = 4 - \frac{6}{n+2}.
\]
This expression is a positive integer if and only if $(n+2) \mid 6$. The positive divisors of $6$ are $1, 2, 3, 6$, yielding:
\begin{align*}
n + 2 = 1 &\implies n = -1 \quad \text{(invalid, since } n \geq 0\text{)}, \\
n + 2 = 2 &\implies n = 0, \\
n + 2 = 3 &\implies n = 1, \\
n + 2 = 6 &\implies n = 4.
\end{align*}
Hence $C_n \mid C_{n+1}$ if and only if $n \in \{0, 1, 4\}$.
\end{proof}

\newpage
\begin{theorem}\label{thm:prime}
If $p = n + 2$ is prime and $p > 5$, then $p \mid C_n$ and
\[
\gcd(C_n, C_{n+1}) = \frac{C_n}{p}.
\]
\end{theorem}

\begin{proof}
Since $C_{n+1}$ is a positive integer, the recurrence relation~\eqref{eq:recurrence} implies that
\[
(n+2) \mid 2(2n+1) \cdot C_n,
\]
i.e., $p \mid 2(2n+1) \cdot C_n$.

By Theorem~\ref{thm:divisibility}, since $p$ is prime and $p > 5$, we have $(n+2) \nmid 6$, and therefore $\frac{2(2n+1)}{n+2} = 4 - \frac{6}{n+2}$ is not an integer. In particular, $p \nmid 2(2n+1)$.

Since $p$ is prime and $p \mid 2(2n+1) \cdot C_n$ but $p \nmid 2(2n+1)$, by Euclid's lemma we conclude that $p \mid C_n$.

Now let $k = C_n / p$. Since $p \mid C_n$, we may write $C_n = p \cdot k$. Substituting into the recurrence relation~\eqref{eq:recurrence}:
\[
C_{n+1} = \frac{2(2n+1)}{p} \cdot p \cdot k = 2(2n+1) \cdot k.
\]
Therefore,
\[
\gcd(C_n, C_{n+1}) = \gcd(p \cdot k,\; 2(2n+1) \cdot k) = k \cdot \gcd(p,\; 2(2n+1)).
\]
Since $p$ is prime and $p \nmid 2(2n+1)$, we have $\gcd(p, 2(2n+1)) = 1$. Hence,
\[
\gcd(C_n, C_{n+1}) = k = \frac{C_n}{p}.
\]
\end{proof}

\section{Computed Data}

We present additional computed values of $\gcd(C_n, C_{n+1})$ to support Theorem~\ref{thm:prime}.

\begin{center}
\begin{tabular}{c|c|c|c|c}
$n$ & $C_n$ & $C_{n+1}$ & $\gcd(C_n, C_{n+1})$ & $n+2$ prime? \\
\hline
0 & 1 & 1 & 1 & Yes \\
1 & 1 & 2 & 1 & Yes \\
2 & 2 & 5 & 1 & No \\
3 & 5 & 14 & 1 & Yes \\
4 & 14 & 42 & 14 & No \\
5 & 42 & 132 & 6 & Yes \\
6 & 132 & 429 & 33 & No \\
7 & 429 & 1430 & 143 & No \\
8 & 1430 & 4862 & 286 & No \\
9 & 4862 & 16796 & 442 & Yes \\
10 & 16796 & 58786 & 8398 & No \\
11 & 58786 & 208012 & 4522 & Yes \\
\end{tabular}
\end{center}

For each row where $n+2$ is prime and $n > 4$, observe that $\gcd(C_n, C_{n+1}) = C_n / (n+2)$.

\section{Conclusion}

We have shown that $C_n \mid C_{n+1}$ only for $n \in \{0, 1, 4\}$, and that when $n + 2$ is a prime greater than $5$, the prime $p = n+2$ necessarily divides $C_n$, yielding a predictable formula for the GCD. A complete characterization of $\gcd(C_n, C_{n+1})$ for composite values of $n+2$ remains an open question for future work.

\begin{thebibliography}{9}
\bibitem{stanley}
R.~P. Stanley, \textit{Catalan Numbers}, Cambridge University Press, 2015.

\bibitem{koshy}
T.~Koshy, \textit{Catalan Numbers with Applications}, Oxford University Press, 2009.
\end{thebibliography}

\end{document}
